\section{Cube comparison and pathwise detection}
\label{cube:section}

The coefficient tilt of Section~\ref{sampling:constrained-section}
produces integer-affine maps $\psi(t)=x+Vt$ on which the constrained
polynomials have small real lifts and integral polynomial parts.
We now show how a correlation detected on many sampled boxes can
be detected on the parent box. The local slices and niltests may
depend on the sampled map.

The first step compares the full sampled cube functional with a
controlled mixture of parent cube functionals. Local detection
then gives an averaged cube norm, and the ambient inverse theorem
gives a parent niltest. Separation converts this detector into
finite approximation in the pathwise testing seminorm.

All these output bounds are fixed before the later sampling scales.
Section~\ref{cube:transfer-section} uses the same ensemble to
recover positive scores, with additional late costs permitted only
in the final patch. Its concluding proof of
Proposition~\ref{sampling:constrained} makes the parameter choices
once for all prescribed requests. Until then, the later scales are
chosen to meet the finite requirements of each estimate.

\subsection{The cube functional to be compared}

For an integer box or parameter slice $B$, let
$\operatorname{Cub}_q(B)$ be the set of tuples
$\mathbf t=(t_0,t_1,\ldots,t_q)$ for which
\[
 t_\omega=t_0+\sum_{i=1}^q\omega_i t_i\in B
 \qquad(\omega\in\{0,1\}^q).
\]
Every cube average uses uniform probability on this set, including
degenerate tuples. For $q=0$ it is the one-site average.
An affine map sends such a tuple to an ambient cube. Nevertheless,
the distribution of that ambient cube need not be uniform: the
coefficient tilt and the number of affine maps producing a given
cube both affect its weight. The following proposition expresses
that weight through functions of the individual vertices.

\begin{proposition}[Cube comparison]
\label{cube:comparison}
Use the coefficient law of Section~\ref{sampling:constrained-section}, with
the numbers of dedicated product blocks and kernel columns chosen
sufficiently large as polynomials of the initial dimensions for
the fixed cube-order and derivative ranges. Assume the hypotheses
of Lemma~\ref{sampling:analytic-removal}.
Fix a pre-cost parameter slice $B$, a bounded cube order $q$, and a pre
accuracy $\varepsilon>0$. The averaged sampled-cube functional
\[
 (P_\omega)_\omega\longmapsto
 \E_\psi\E_{\mathbf t\in\operatorname{Cub}_q(B)}
                  \prod_\omega P_\omega(\psi(t_\omega))
\]
is, uniformly for arbitrary $|P_\omega|\le1$, within
$\varepsilon$ of a mixture of ambient cube functionals
\[
 \sum_\alpha c_\alpha
 \E_{\mathbf u\in\operatorname{Cub}_q(U)}
       \prod_\omega P_\omega(u_\omega)J_{\alpha,\omega}(u_\omega).
\]
The $J_{\alpha,\omega}$ are allowed site twists of pre complexity,
and $\sum_\alpha|c_\alpha|$ is pre-bounded. Mixtures can be
replaced by finite sums to any prescribed accuracy. The bounds do
not use $\sigma,L,R$, or the heights of the subspaces $W_h$.
\end{proposition}

The functions $P_\omega$ may be arbitrary bounded inputs. Thus
equidistribution of the coefficient arrays alone does not prove
the proposition. We use Lemma~\ref{sampling:analytic-removal}
with these vertex functions retained: only coefficient modes
determined by evaluation at the cube vertices survive. Counting
the spatial fibers then gives the weight of the resulting ambient
cube. The remaining problem is to describe the distribution of
the evaluated small lifts.

There are two kinds of lift coordinates. Coordinates parallel to
$W_h$ have real-valued laws, while projected integer coordinates
have lattice laws at scales $K_i$. In both cases the principal
polynomial is a sum of products on disjoint parameter blocks.
Those products supply a smooth distribution of the evaluated
values. For the real coordinates we must also justify removing
the small nonprincipal coefficients. For the lattice coordinates
we retain the congruence factors as well as the smooth dependence
on $z_i/K_i$.

We first prove the product-block estimate, then count the spatial
and coefficient fibers, and finally compare the real image laws.
These three estimates are assembled in the proof of
Proposition~\ref{cube:comparison}. The lift estimates act on the
fixed number of outputs associated with one lift axis. The number
of axes and input parameters may grow. Their counts enter the
bounds polynomially; caps and inverse precisions enter through
their logarithms.

\subsection{Evaluating the product blocks}

Fix a degree $h\ge1$ and a cube order $q\ge0$, both in the
prescribed fixed range. Evaluation of a degree-$h$ polynomial
$P$ on a parameter cube has the form
\[
 P\left(t_0+\sum_{i=1}^q\omega_it_i\right)
 =\sum_{\substack{S\subseteq[q]\\|S|\le h}}
       P_S(t_0,\ldots,t_q)\prod_{i\in S}\omega_i,
 \qquad\omega\in\{0,1\}^q.
\]
Powers of each $\omega_i$ are reduced by $\omega_i^2=\omega_i$.
The vector $(P_S)_S$, called the Boolean jet, determines all the
vertex values. It has
\begin{equation}
 j(h,q)=\sum_{a=0}^{\min(h,q)}\binom qa
 \label{cube:jet-dimension}
\end{equation}
coordinates. This dimension is independent of the number of
parameter columns and lift axes. At $q=0$ the jet consists of a
single value; the same estimate will therefore control the
one-site weights used in score recovery.

For a sum of independent product blocks, the Fourier coefficients
of its jet law multiply. We will show that a single block can
have substantial Fourier bias only at a bounded rational
frequency plus a bounded frequency in normalized coordinates.
Taking a sufficiently large structural number of blocks then
gives an absolutely summable Fourier expansion. The required
block count is fixed before the eventual truncation accuracy.
The next two elementary estimates establish the single-block
bias assertion.

\begin{lemma}[A many-multiples estimate]
\label{cube:many-multiples}
There are absolute constants $C>1$ and $0<c<1$ with the following property.
Let $0<\eta\le1$, $0<\delta<c\eta$, and $N>C\eta^{-2}$.
If at least $\eta N$ integers $v\in[-N,N]$ satisfy
$\|v\theta\|_{\R/\Z}\le\delta$, then there is an integer
$1\le d\le C\eta^{-1}$ such that
\[
                  \|d\theta\|_{\R/\Z}
                    \le \frac{C\delta}{\eta^2N}.
\]
The same assertion, with changed absolute constants, holds for any
interval of length comparable to $N$.
\end{lemma}

\begin{proof}
Two successive good integers have positive difference
$d\le C/\eta$, so $\gamma=\|d\theta\|\le2\delta$.
There is nothing to prove if $\gamma=0$.  Split $[-N,N]$ into
progressions of common difference $d$.  On each progression, the
successive points on the circle are spaced by $\gamma$, in either
orientation.  A progression of length $M$ therefore visits an arc of
length $2\delta$ at most
\[
       C\bigl((\delta+\gamma)M+1+\delta/\gamma\bigr)
\]
times: split the lifted interval of length $M\gamma$ at the integers,
and count at most $1+2\delta/\gamma$ points in each resulting visit.
Summing over the at most $d$ progressions gives
\[
 \eta N\le C\bigl((\delta+\gamma)N+d+d\delta/\gamma\bigr).
\]
For a sufficiently small absolute $c$ and a sufficiently large
constant in the lower bound for $N$, the first two terms are at most
$\eta N/2$.  Thus
$\gamma\le Cd\delta/(\eta N)\le C\delta/(\eta^2N)$.
\end{proof}

\begin{lemma}[Dense products and multilinear bias]
\label{cube:multilinear-bias}
For every fixed $e\ge1$ there are constants $C_e,c_e>0$ with the
following properties.

First, suppose a set of at least an $\eta$ fraction of the integer
tuples $v_i\in[-N_i,N_i]$, $1\le i\le e$, satisfies
$\|\theta v_1\cdots v_e\|\le\delta$.  If
\[
 \min N_i\ge C_e\eta^{-C_e},\qquad
                  \delta\le c_e\eta^{C_e},
\]
then some integer $1\le d\le C_e\eta^{-C_e}$ satisfies
\[
             \|d\theta\|\le
                   \frac{C_e\eta^{-C_e}\delta}{N_1\cdots N_e}.
\]
Second, let $F(x_1,\ldots,x_e)$ be a real polynomial of degree at
most one in each variable, with coefficient $\theta$ on
$x_1\cdots x_e$.  Let the independent integer variables $x_i$
range uniformly over intervals of lengths $N_i$.  If
\[
 |\E e(F(x))|\ge\zeta>0,
 \qquad\min N_i\ge C_e\zeta^{-C_e},
\]
then $\theta$ is within
$C_e\zeta^{-C_e}/(N_1\cdots N_e)$ of a rational with denominator
at most $C_e\zeta^{-C_e}$.
\end{lemma}

\begin{proof}
For the first assertion use induction on $e$, the case $e=1$ being
Lemma~\ref{cube:many-multiples}.  A fraction at least $c\eta$ of
the outer tuples $(v_2,\ldots,v_e)$ have a fiber with at least
$c\eta N_1$ good values of $v_1$.  On each such fiber the preceding
Lemma gives an integer $d_1\le C/\eta$ such that
\[
 \|d_1\theta v_2\cdots v_e\|
                    \le \frac{C\delta}{\eta^2N_1}.
\]
One value of $d_1$ works on a fraction at least $c\eta^2$ of the
outer tuples.  Apply the induction hypothesis with density
$c\eta^2$, tolerance $C\delta/(\eta^2N_1)$, and coefficient
$d_1\theta$.  Increasing $C_e$ and decreasing $c_e$ in terms of
$C_{e-1},c_{e-1}$ makes all of its hypotheses hold.  Its denominator
multiplied by $d_1$, and its error multiplied by
$C/\eta^2$, have the asserted form.  For example a recurrence with
$C_e$ greater than $4C_{e-1}+10$, together with sufficiently small
$c_e$, suffices.  The number of iterations is fixed.

For the second assertion the case $e=1$ is the geometric-sum bound.
For larger $e$, apply box Cauchy--Schwarz in the first $e-1$
variables, using independent pairs $x_i,x_i'$.  All terms missing
one of these variables disappear after differencing.  The only term
linear in $x_e$ that remains has coefficient
$\theta\prod_{i<e}(x_i-x_i')$; the remaining terms are constant
in $x_e$ and disappear on taking its absolute mean.  Hence, with
$a=\zeta^{2^{e-1}}$,
\[
 a\le\E_{x_i,x_i',\ i<e}
 \min\left(1,\frac{2}{N_e
         \|\theta\prod_{i<e}(x_i-x_i')\|}\right).
\]
For a fraction at least $a/2$ of these pairs the norm in the
denominator is at most $4/(aN_e)$.  Each difference value has
probability at most $1/N_i$.  Thus the corresponding difference
tuples occupy a fraction at least $c_ea$ of the full product of
integer intervals $[-N_i,N_i]$, $i<e$.  Apply the first assertion
with $\eta=c_ea$ and $\delta=4/(aN_e)$.  The prescribed sufficiently
large power lower bound on all the $N_i$ ensures its size and
tolerance hypotheses.  Divide the resulting estimate for
$\|d\theta\|$ by $d$; increasing $C_e$ one last time gives the
claimed rational approximation.
\end{proof}

\paragraph{The laws to be smoothed.}
Fix $h,q$ and write $j=j(h,q)$. In one projection direction
the principal polynomial is a sum of $J_0$ products, each using
its own $h$ parameter coordinates. Its constant term is fixed
or is an independent integer shift. There are two discrete laws.
On an inactive axis, each product has coefficient $1$ and its
parameter sides are comparable to $cK^{1/h}$. On a moderate
axis, its parameter sides are $L$ and its coefficient has an
independent normalized smooth integer law on an interval of
length comparable to $c'K/L^h$. The constants $c,c'$ have pre
bounds. A zero principal polynomial is allowed only when $K$
is pre-bounded; that case will be enumerated.

Cube parameters on different axes are independent. On a single
axis their normalized domain is the polytope on which all cube
vertices belong to its designated interval. We allow
coordinatewise slices and conditioning modulo pre-bounded
moduli. These restrictions preserve independence between axes
and between the dedicated blocks. The moderate block
coefficients are also independent of all parameter tuples.

For the continuous law, use real normalized cube parameters on
these polytopes. Each principal coefficient is either fixed
with magnitude in $[B^{-1},B]$ or sampled independently with the
specified normalized smooth density supported in that range.
The bound $B$ also includes all inverse normalized volumes,
coefficient widths, smooth-density bounds and residue moduli.

An allowed slice of relative size at least $\delta$ has stride
at most $2/\delta$ on a sufficiently long axis. Its normalized
scalar cube polytope has dimension $q+1$, at most $2^{q+1}$
faces and volume at least $c_q\delta^{q+1}$. Its density is
therefore at most $C_q\delta^{-(q+1)}$. Conditioning modulo $M$
multiplies the stride by at most $M$. We include these pre
losses in $B$. If an inactive axis is too short for these
comparisons, its $K$ is pre-bounded and is enumerated instead.

We use a finite Fourier expansion for the discrete law. The
normalized jet support lies in a fixed box. Choose a fixed
positive integer $R_{\rm tor}$ so this box is contained in a
fundamental box for $\R^j/R_{\rm tor}\Z^j$. If $Y$ is the
integer jet, set
\[
 D_K(z)=K^j\Pr(Y=z),
 \qquad z\in(\Z/(R_{\rm tor}K)\Z)^j,
\]
using the genuine integer representatives on its support.
This is density relative to grid volume $K^{-j}$. Fourier
inversion reads
\[
 D_K(z)=R_{\rm tor}^{-j}\sum_\xi
 \E e\left(-\frac{\xi\cdot Y}{R_{\rm tor}K}\right)
                 e\left(\frac{\xi\cdot z}{R_{\rm tor}K}\right).
\]
The sum is over the displayed finite group. The corresponding
continuous convention is Fourier series on
$\R^j/R_{\rm tor}\Z^j$, with a smooth cutoff outside the
small support.

\begin{lemma}[Uniform product-block Fourier estimate]
\label{cube:product-block}
For the preceding laws, choose the number $J_0$ of disjoint
blocks sufficiently large as a polynomial in the initial
dimensions. In the moderate case use the sampling scale range
\[
 K\le L^{C_s},\qquad K/L^h\ge A^h,
 \qquad \log L\le C_0d\log A,
 \qquad d=1+\sum_h d_h,
\]
where $C_s,C_0$ depend only on the fixed degree range.

The Fourier coefficients of $D_K$ have a pre-bounded absolute
sum, uniformly in $K,L$ and in the permitted slices and residue
conditions. For every pre error $\varepsilon>0$, retaining
pre-boundedly many characters gives pointwise error at most
$\varepsilon$. Their coefficient sum remains pre-bounded,
and every retained character can be written as
\begin{equation}
 e\bigl((a/d'+\xi/K)\cdot z\bigr),
 \qquad d',|\xi|\le\mathcal B(b_{\rm pre}),
 \label{cube:major-character}
\end{equation}
with $a\in\Z^j$, $\xi\in\R^j$, and a positive integer denominator $d'$.
Thus one factor is a bounded-modulus character and the other
is smooth in $z/K$.

The continuous jet law has a density with pre-bounded
derivatives through any fixed order chosen before $J_0$.
The choice of $J_0$ uses this derivative order and the initial
dimensions, but not the eventual truncation accuracy. Fixed
constant shifts preserve all these bounds.
\end{lemma}

\begin{proof}
We first locate the frequencies at which one block can have
large bias. We then use the scale range to make this estimate
uniform down to a power of $K^{-1}$, and multiply the estimates
over the independent blocks.

Testing one product block against a jet character gives a
multilinear phase in $h$ vectors of $q+1$ cube coordinates.
A sampled product coefficient supplies one further scalar
variable. Every jet frequency is an entry of this multilinear
form, up to the fixed principal coefficient: specialize its
affine factors to $1$ or to selected Boolean coordinates to
isolate a squarefree monomial of degree at most $h$.

Suppose the bias has magnitude at least $\zeta$. Subdivide
each normalized cube polytope into boxes of mesh
$(\zeta/B)^{C_{h,q}}$. A boundary strip of width $\delta$
has volume at most $C_{h,q}\delta$, since its dimension and
number of faces are fixed. The corresponding discrete count
has additional error at most $C_{h,q}\sum_iN_i^{-1}$.
After division by the domain volume, which is at least
$B^{-1}$, these boundary errors are small compared with
$\zeta$. The smooth weight varies by at most
$C_{h,q}B\delta$ on a box. Consequently, if
\[
                    \min_iN_i\ge(B/\zeta)^{C_{h,q}},
\]
one retained box has unweighted bias at least $\zeta/4$.
Its side lengths are at least a fixed power of $\zeta/B$
times the original lengths. Residue conditioning gives the
same conclusion after reindexing by its pre-bounded stride;
its counting error is controlled by the reciprocal reindexed
side lengths.

Choose a tensor entry and freeze all other scalar coordinates
at values retaining that bias. The chosen full multilinear
coefficient is unchanged, because every monomial uses at most
one scalar coordinate from each vector group. Apply
Lemma~\ref{cube:multilinear-bias} to the remaining $e=h$ or
$h+1$ variables. It shows that this entry is within
\[
              \frac{(B/\zeta)^{C_{h,q}}}{N_1\cdots N_e}
\]
of a rational with denominator at most
$(B/\zeta)^{C_{h,q}}$. Different entries may use different
favorable frozen coordinates; their full multilinear
coefficients do not depend on those choices. The product
$N_1\cdots N_e$ is comparable to $K$ within pre factors.
There are only $j$ entries to control, so their denominators
may be combined at the same type of pre cost.

We next check the side-length requirement uniformly in the
sampling scale. Enlarge $C_0$ so that $C_0\ge s$. On a
moderate axis,
\[
 L\ge K^{1/C_s},\qquad
 K/L^h\ge A^h\ge L^{h/(C_0d)}
                 \ge K^{h/(C_sC_0d)}.
\]
Thus every scalar side, including the coefficient interval,
is at least $B^{-1}K^\kappa$, with
$\kappa=(C_sC_0d)^{-1}$. Inactive sides satisfy the same
weaker bound because they have size at least
$B^{-1}K^{1/h}$. Let $C$ be the exponent required by the
single-block argument and choose
\[
                   P_0\ge2C/\kappa,
                   \qquad B'\ge B^{1+1/C}.
\]
At $\zeta=B'K^{-1/P_0}$ one has
\[
                 (B/\zeta)^C\le B^{-1}K^{\kappa/2},
\]
so the side-length requirement holds. It also holds for
larger $\zeta$. The exponent $P_0$ depends polynomially on
the initial dimensions and not on the requested accuracy.
If this lower bias cutoff exceeds a desired retained level
$\zeta_0$, then $K\le(B'/\zeta_0)^{P_0}$ is pre-bounded;
the whole grid can be enumerated.

At a bias level $\zeta$, the possible major frequencies
number at most $(B/\zeta)^{C_{h,q}}$. Indeed the denominator
and numerator ranges have those bounds, and each interval of
width $(B/\zeta)^{C_{h,q}}/K$ contains only that many points
of the Fourier grid. The number $j$ of coordinates is fixed.
Choose common nested major sets at dyadic bias levels, using
the uniform estimate for every block. Outside the major set
at level $\zeta$, each single-block Fourier coefficient has
modulus at most $\zeta$. Independence gives modulus at most
$\zeta^{J_0}$ for the sum of the blocks.

Summing over the levels above the cutoff gives a convergent
geometric sum when $J_0$ is sufficiently large. Below the
cutoff, the trivial number $O(K^j)$ of grid frequencies gives
the bound
\[
                    O(K^j)(B'K^{-1/P_0})^{J_0}.
\]
If the cutoff is below $\zeta_0$, this is at most a fixed
multiple of
\[
                 (B')^{P_0j}\zeta_0^{J_0-P_0j}.
\]
Take $J_0>2P_0j+2C_{h,q}$. The two discarded sums become
as small as requested by decreasing $\zeta_0$ at pre cost.
Fixing $\zeta_0$ once at a constant level gives the absolute
Fourier cap independently of the final accuracy. The retained
major frequencies have precisely the form
\eqref{cube:major-character}.

For the continuous law, suppose a single-block Fourier
coefficient at frequency $\xi$ has magnitude at least
$\beta>0$. Approximate the normalized real variables by
grids of mesh $1/m$. For all sufficiently large $m$ the
discrete bias is at least $\beta/2$. A selected full
multilinear coefficient becomes $t(\xi)/m^e$, where
$t(\xi)$ is its original linear functional of the jet
frequency. The discrete estimate supplies
\[
 1\le d'\le H,\qquad
 \left|\frac{t(\xi)}{m^e}-\frac{a}{d'}\right|
       \le\frac{H}{m^e},\qquad
 H=(CB/\beta)^{C_{h,q}}.
\]
The bound for $H$ is independent of $m$. Taking
$m^e>H(|t(\xi)|+H)$ forces $|a/d'|<1/H$, hence $a/d'=0$
and $|t(\xi)|\le H$. The tensor entries include every jet
frequency, up to a pre-bounded nonzero principal coefficient.
It follows that $|\xi|\le BH$, giving single-block decay
$B'(1+|\xi|)^{-c_{h,q}}$.

Independent blocks multiply this decay. For the derivative
order $r$ fixed initially, choose $c_{h,q}J_0>j+r$.
Fourier inversion with integrable weight $1+|\xi|^r$
gives the claimed derivative bounds. Fixed shifts multiply
Fourier coefficients by phases and leave every estimate
unchanged.
\end{proof}

\subsection{Counting integer fibers}

The spatial affine map sends integer inputs to the root and
differences of the sampled cube. Its pushforward consists of a
smooth real density multiplied by an integer-image mask. The next
lemma identifies both factors. Applied one ambient coordinate or
lift axis at a time, its mesh loss has an exponent depending on
the fixed output rank rather than the growing number of input
coordinates.

\begin{lemma}[Smooth counting on a linear fiber]
\label{cube:fiber}
Let $A:\Z^e\to\Z^j$ have real rank $j$, integer entries bounded by
$CL^h$, and image of index $g$ in $\Z^j$.  Let positive diagonal
input and output scales be $S_\nu$ and $P_i$.  Assume the rescaled
matrix $\widetilde A_{i\nu}=P_i^{-1}A_{i\nu}S_\nu$ has bounded
entries and has a $j$-minor with bounded inverse.  Let $\varphi$ be a
$C^1$ function supported in a fixed bounded set in normalized input
coordinates.  Give $z\in\Z^e$ mass
$(\prod S_\nu)^{-1}\varphi((z_\nu/S_\nu)_\nu)$.

At an output point $v\in A\Z^e$, its mass, multiplied by
$\prod P_i$, differs from $g$ times the continuous normalized
pushforward density by at most
\begin{equation}
 C_\varphi\,\frac{L^{hj}}{\min_\nu S_\nu}.
 \label{cube:fiber-error}
\end{equation}
Here $C_\varphi$ depends on $e,j,C$, the support and $C^1$ bounds,
and the normalized minor and index bounds.  It is independent of the
scales.  The assertion is used only when the displayed mesh is small.
Off $A\Z^e$ the discrete mass is zero.
\end{lemma}

\begin{proof}
Write $J(B)=\sqrt{\det(BB^{\mathsf T})}$ for a full-row-rank matrix,
and let $S,P$ be the diagonal scale matrices.  Choose a nonsingular
$j$-column pivot $A_0$.  For each remaining column $a_\nu$, the
vector with pivot part $-\operatorname{adj}(A_0)a_\nu$ and free
coordinate $\det A_0$ is an integer kernel vector.  These $e-j$
vectors are independent and their coordinates are bounded by
$j!C^jL^{hj}$.  Their fundamental parallelepiped bounds the covering
radius of their span, and hence that of the full integer kernel, by
$C_{e,j,C}L^{hj}$.  Voronoi cells of the full kernel have diameter
at most twice this bound.  Thus the number of input variables enters
the constant, not the exponent of $L$.

If $\Lambda=\ker(A)\cap\Z^e$, then
\begin{equation}
 \covol(\Lambda)=\frac{J(A)}{g}.
 \label{cube:kernel-covolume}
\end{equation}
Indeed $\Lambda$ is primitive.  Extend its basis by lifts of a basis
of $A\Z^e$ to obtain a basis of $\Z^e$.  Orthogonal projection onto
$\ker(A)^\perp$, on which $A$ has volume factor $J(A)$, gives
$1=\covol(\Lambda)g/J(A)$.

The restriction of $S$ from $\ker(\widetilde A)$ to $\ker(A)$ has
Jacobian $\det(S)J(A)/(\det(P)J(\widetilde A))$, by coarea applied
to $AS=P\widetilde A$.  Consequently the normalized kernel lattice
$\Lambda'=S^{-1}\Lambda$ satisfies the exact identity
\[
 \covol(\Lambda')=
 \frac{\det(P)J(\widetilde A)}{g\det(S)}.
\]
Transform the original kernel cells by $S^{-1}$.  They tile each
attained normalized fiber and have diameter at most
$r=C_{e,j,C}L^{hj}/\min S_\nu$.

Suppose $\operatorname{supp}\varphi\subset[-R_0,R_0]^e$, $r\le1$,
and the inverse norm of the normalized pivot is at most $K_0$.
Solving for its $j$ coordinates gives, uniformly in the output $y$,
\[
 \frac{\mathcal H^{e-j}
  (\{\widetilde Ax=y\}\cap[-R_0-1,R_0+1]^e)}
      {J(\widetilde A)}
 \le K_0^j(2R_0+2)^{e-j}.
\]
The continuous normalized output density is
$f(y)=J(\widetilde A)^{-1}\int_{\widetilde Ax=y}\varphi\,
d\mathcal H^{e-j}$.  The covolume identity shows that the
normalized discrete mass is a cellwise Riemann sum for $g f(y)$.
The oscillation on each cell is at most
$\|\nabla\varphi\|_\infty r$; summing only cells meeting the
support therefore bounds its error by
\[
 g\|\nabla\varphi\|_\infty r K_0^j(2R_0+2)^{e-j}.
\]
This proves~\eqref{cube:fiber-error}, with explicit dimensional
growth.  For a continuous input mass one, ordinary full-lattice
cells also give a normalization error
$O_{e,R_0}(\|\varphi\|_{C^1}/\min S_\nu)$.  When this is small,
division by the discrete total mass preserves the estimate, using
the same uniform bound for $f$.
\end{proof}

\begin{lemma}[Fixed-dimensional polynomial sublevels]
\label{cube:polynomial-sublevel}
Let $N,d\ge1$.  If a real polynomial $Q$ in $N$ variables has degree
at most $d$ in each variable and has a coefficient of magnitude at
least $c>0$, then, for $0<u\le c$,
\[
 \big|\{x\in[-1,1]^N:|Q(x)|\le u\}\big|
 \le C_{N,d}(u/c)^{1/(Nd)}.
\]
The same conclusion holds on any fixed containing box, with changed
constants.  A probability density bounded by $H$ multiplies the
right side by $H$.
\end{lemma}
\begin{proof}
In one variable, if the sublevel set has length $\lambda>0$, choose
$d+1$ points in it separated by at least $\lambda/(4d+4)$.
Lagrange interpolation at these points bounds every coefficient by
$C_d u\lambda^{-d}$.  This gives the assertion for $N=1$.
For induction, expand in the last variable and choose a coefficient
polynomial having a coefficient of magnitude at least $c$.
Discard the base points where its magnitude is at most $\tau$.
The induction hypothesis and the one-variable estimate bound the
total sublevel volume by
\[
 C_{N,d}\left((\tau/c)^{1/((N-1)d)}+(u/\tau)^{1/d}\right).
\]
Set $\tau/c=(u/c)^{(N-1)/N}$.  Rescaling a fixed box changes
coefficient norms by fixed factors and proves the remaining claims.
\end{proof}

\subsection{Kernel tuples and finite evaluation covers}

To compare a sampled cube with a parent cube, we will solve for
spatial variables using the kernel columns. A retained parameter
tuple must permit this operation over both the reals and the
integers. A real minor controls the smooth spatial density; the
integer cokernel determines a finite cover on which coefficient
characters can be read as functions of the evaluated jet.

\begin{lemma}[Kernel tuples and evaluation covers]
\label{cube:kernel-covers}
Let the cube order $q$ and the maximum polynomial degree $s$
belong to the fixed ranges, and use the independent kernel-column
laws with sides $L$ and the prescribed coordinatewise slice and
residue restrictions. Choose a structural number $J>q(q+2)$
of kernel columns when $q\ge1$.
For every pre error $\eta>0$, there is a pre bound $B$ such that,
for all sufficiently large later $L$, all but probability $\eta$
of the parameter cube tuples have the following properties.

Let $D$ be the $q\times J$ matrix of their differences on the
kernel columns. It has a real $q$-minor of magnitude at least
$L^q/B$, and the finite cokernel
$\Z^q/D\Z^J$ has exponent $a\le B$.
For such a tuple, let $E_{h,t}$ be the integer matrix taking
ordinary coefficients of degree-at-most-$h$ parameter
polynomials to their Boolean jets, and let $e_h$ be the number
of these parameter monomials. For every $h\le s$,
this map has a rational right inverse with denominators
dividing $a^h$.

Consequently $M=a^s$, or its product with a previously prescribed
pre-bounded modulus, satisfies
\[
                  E_{h,t}\Z^{e_h}\supset M\Z^{j(h,q)}.
\]
Membership in the image lattice is determined modulo $M$, and
its index is
\[
       [\Z^{j(h,q)}:E_{h,t}\Z^{e_h}]
        =\frac{M^{j(h,q)}}{|\operatorname{im}(E_{h,t}\bmod M)|}.
\]
Both its membership mask and its index are determined by the
parameter tuple modulo $M$.

For $q=0$ every tuple has these conclusions: the difference
matrix has zero rows, its empty minor is $1$, and its cokernel
has exponent $1$. Constant polynomials give an integral right
inverse for the one-site evaluation map. One may take $M=1$,
or the previously prescribed modulus, without discarding a set.
\end{lemma}

\begin{proof}
Assume $q\ge1$. Choose any fixed $q$ kernel columns. Their
normalized difference determinant is a polynomial of degree
$q$ with a coefficient of magnitude one. The normalized scalar
cube domains have pre-bounded densities on fixed containing
boxes. Lemma~\ref{cube:polynomial-sublevel} therefore bounds
the probability that the determinant has magnitude at most
$2/B$ by a pre constant times a fixed power of $B^{-1}$.
For the lattice parameter law, majorize the bad event by a
smooth function vanishing beyond $2/B$. Its Lipschitz norm
has pre cost, so grid approximation adds an error tending to
zero with $L$. This proves the real-minor assertion.

On the full-rank event, the cokernel exponent divides a
nonzero $q$-minor, and is therefore at most $C_qL^q$.
A character of exact order $a$ on $\Z^q$ has some coordinate
of order at least $a^{1/q}$. After fixing the other entries
of one kernel column, annihilation by this character requires
that selected difference entry to lie in at most one residue
class of this order. Counting in a containing box and dividing
by the scalar cube-domain mass bounds its probability by
\[
                       B'(a^{-1/q}+L^{-1}).
\]
The prescribed strides and residue conditions reduce the
effective order by at most their pre-bounded modulus. Absorb
this loss into $B'$. For $a\le C_qL^q$ the displayed bound
is at most $B''a^{-1/q}$.

The kernel columns are independent before imposing any
full-rank condition. If the matrix has full rank and cokernel
exponent greater than $B$, a character of some exact order
$B<a\le C_qL^q$ annihilates every column. There are at most
$a^q$ such characters at order $a$. Thus
\[
 \Pr(\text{full rank and exponent}>B)
 \le (B'')^J\sum_{a>B}a^{q-J/q}
 \le C_{q,J}(B'')^J B^{q+1-J/q}.
\]
The structural choice $J>q(q+2)$ makes this tail decrease
as a fixed power of $B$. Choose $B$ for both exceptional
sets and then choose $L$ for the grid errors. The number of
columns has not been increased in response to $\eta$.

Fix a retained tuple, and write $t_K$ for its kernel
coordinates. For each $1\le r\le q$, choose
$n_r\in\Z^J$ with $Dn_r=ae_r$. The rational affine
polynomial
\[
                    \ell_r(t)=
                    n_r\cdot(t_K-t_{0,K})/a
\]
equals $\omega_r$ at the vertex $t_\omega$. Products of at
most $h$ of these polynomials, together with the constant,
give a right inverse for the Boolean-jet map with
denominators dividing $a^h$. This proves the image inclusion
with $M=a^s$; multiplying $M$ by another prescribed modulus
preserves it. Reduction modulo $M$ now gives the asserted
membership and index formulas.

If numerator bounds are needed, choose a nonzero pivot minor
of $D$, reduce nonpivot coordinates of the solutions modulo
its determinant, and solve for the pivot coordinates using
its adjugate. This bounds the interpolation numerators by
powers of $L$. The cover assertion itself needs only the
bounded denominator.
\end{proof}

The cover may depend on the retained kernel tuple. No common
multiple over all possible exponents $a\le B$ is taken.
This is sufficient because the subsequent cube comparison
first fixes that tuple and works on its individual cover.

\subsection{Stability of nonsingular polynomial images}

The real output here consists of the continuous lift axes and the
enormous projection axes after their lattice laws have been replaced
by continuous laws. We condition on all variables generating inactive
and moderate grid outputs. On each remaining axis, independent
principal product blocks supply a nonsingular minor. The following
statement then allows us to remove the small nonprincipal coefficients
jointly in the real outputs, uniformly in that conditioning. It does
not assert total variation stability for perturbations of a discrete
law. We use the total variation
norm $\|\mu\|_{\mathrm{TV}}=\sup_{|f|\leq 1}|\int f\,d\mu|$; thus differences
of probability measures have norm at most $2$.

\begin{lemma}[Perturbing a polynomial image]\label{cube:image-perturbation}
Let $\Omega\subset\mathbb R^a$ be bounded, let $\rho\,dx$ be a probability
measure on $\Omega$, and let $F,G:\Omega\to\mathbb R^b$ be $C^2$ maps on a
neighbourhood of $\overline\Omega$. Suppose that there is a function
$\chi\in C_c^1(\Omega)$ with $0\leq\chi\leq1$ and
\[
 \int_\Omega(1-\chi)\rho\leq\eta,
 \qquad w:=\chi\rho\in C_c^1(\Omega).
\]
Fix $b$ input coordinates $i_1,\ldots,i_b$, and write
\[
 A_t(x)=\big(\partial_{i_1}F_t(x),\ldots,
                  \partial_{i_b}F_t(x)\big),\qquad
 F_t=F+t(G-F),\quad 0\leq t\leq1.
\]
Assume $A_t$ is invertible on a neighbourhood of $\operatorname{supp}w$.
For $1\leq r\leq b$, let $V_{t,r}$ be the input vector field supported in
these $b$ coordinate directions whose selected components are
$A_t^{-1}e_r$. Suppose
\[
 \sup_{t,r}\int_\Omega
       \big|\operatorname{div}(wV_{t,r})\big|\,dx\leq B.
\]
If $\|F-G\|_{L^\infty(\operatorname{supp}w),\ell^\infty}\leq\varepsilon$,
then, for every $\delta>0$,
\[
 \|F_*(\rho\,dx)-G_*(\rho\,dx)\|_{\mathrm{TV}}
 \leq 2\eta+2bB\delta+C b\varepsilon/\delta,
\]
where $C$ is an absolute constant. In particular, when $B\geq1$ and
$0<\varepsilon\leq1$, the right-hand side is at most
$2\eta+C' b\sqrt{B\varepsilon}$.
\end{lemma}

\begin{proof}
Put $\mu_t=(F_t)_*(w\,dx)$.  For a smooth compactly supported test
function $\varphi$, the chain rule and integration by parts give
\begin{align*}
 \int_{\mathbb R^b}\partial_r\varphi\,d\mu_t
 &=\int_\Omega w\,V_{t,r}\cdot\nabla(\varphi\circ F_t)\,dx\\
 &=-\int_\Omega\varphi(F_t(x))
                   \operatorname{div}(wV_{t,r})(x)\,dx.
\end{align*}
Consequently each distributional derivative of $\mu_t$ is a finite signed
measure of total variation at most $B$. Integrating the corresponding
identity along a coordinate translation and then summing over coordinates
shows
\[
 \|\tau_z\mu_t-\mu_t\|_{\mathrm{TV}}\leq B\|z\|_1.
\]
One may first test this identity on smooth functions; regularity of finite
Borel measures then gives the stated total variation bound.

Choose a nonnegative one-dimensional smooth mollifier $\vartheta$, of
integral $1$ and support in $[-1,1]$, and let
$\vartheta_\delta^{\otimes b}$ be its product at scale $\delta$. Averaging
the translation estimate gives
\[
 \|\mu_t*\vartheta_\delta^{\otimes b}-\mu_t\|_{\mathrm{TV}}
 \leq bB\delta.
\]
In particular $\mu_t$ is absolutely continuous: it is a total variation
limit of the absolutely continuous measures on the left. Its density belongs to the space $BV$ of functions whose distributional
first derivatives are finite measures, with variation at most $B$ in each
coordinate.

The measure $w\,dx$, of mass at most $1$, couples $\mu_0$ and $\mu_1$.
The $L^1$ norm of each partial derivative of the product mollifier is at
most $C/\delta$. Hence this coupling gives
\[
 \|\mu_0*\vartheta_\delta^{\otimes b}
          -\mu_1*\vartheta_\delta^{\otimes b}\|_{\mathrm{TV}}
 \leq Cb\varepsilon/\delta.
\]
The triangle inequality proves the displayed estimate for the cutoff
measures. The discarded measures each have mass at most $\eta$, giving
$2\eta$ for the original measures. Take
$\delta=\sqrt{\varepsilon/B}$ for the last assertion.
\end{proof}

\begin{remark}[Checking the hypotheses and their costs]\label{cube:perturbation-cost}
It is enough that, on the support of $w$, the selected minor of $DF$ has
inverse norm at most $K$, that $F,G$ have $C^2$ norm at most $H$, and that
$\|DF-DG\|\leq(2K)^{-1}$. A Neumann series gives
$\|A_t^{-1}\|\leq2K$, uniformly in $t$, and differentiating an inverse
matrix gives
$\|\partial_i A_t^{-1}\|\leq4K^2\|\partial_i A_t\|$.
The product rule therefore bounds the required $B$ by a polynomial in
$a,b,K,H,\|w\|_{W^{1,1}}$ (with a harmless choice of matrix norms).
For example, taking $H$ to bound each selected second-derivative
matrix in operator norm, one may use
\[
 B\le 2K\|\nabla w\|_{L^1}
          +4bK^2H\|w\|_{L^1}.
\]
All these estimates concern the good-domain cutoff only; no density bound
is required on the discarded set.

There is also a sharper estimate if $\|F-G\|_{C^1}\leq\varepsilon$.
Let $U_t$ be the input vector field with selected components
$A_t^{-1}(G-F)$. Differentiating in $t$ and integrating by parts gives
\[
 \frac{d}{dt}\int\varphi(F_t(x))w(x)\,dx
 =-\int\varphi(F_t(x))\operatorname{div}(wU_t)(x)\,dx.
\]
The same inverse-matrix bounds show
$\int|\operatorname{div}(wU_t)|\leq B'\varepsilon$, where $B'$ has the
same type of polynomial dependence. Integration in $t$ yields the
stronger bound $2\eta+B'\varepsilon$. The mollification proof above is
useful because it needs only a uniform displacement bound once the
first-derivative estimates for the two image measures have been supplied.
\end{remark}

\paragraph{The required minors for product blocks.}
Fix integers $h\geq1$, $q\geq0$, and put
\[
 \mathcal S=\{S\subseteq\{1,\ldots,q\}:|S|\leq h\},
 \qquad j=|\mathcal S|.
\]
For an affine factor write
$\ell(\omega)=a_0+\sum_{r=1}^q a_r\omega_r$.
Let $\mathcal J$ send a polynomial on the Boolean cube to its coefficients
in the squarefree monomial basis $\omega_S$. Consider the map
\[
 F_{\mathrm{ax}}(a)
   =\mathcal J\left(\sum_{B=1}^{J}
               c_B\prod_{v=1}^h\ell_{B,v}(\omega)\right),
 \qquad 0<c_0\leq |c_B|\leq c_1,
\]
where different blocks use disjoint input coordinates and $J\geq j$.
Allocate a different block $B(S)$ to each $S\in\mathcal S$.
If $|S|<h$, differentiate its first factor with respect to its constant
coefficient; specialize $|S|$ of the remaining factors to the different
variables $\omega_r$, $r\in S$, and all other factors to $1$.
If $|S|=h$, differentiate its first factor with respect to a coefficient
$a_r$ with $r\in S$; specialize the other factors to the remaining
$h-1$ variables. In either case the resulting Jacobian column is exactly
$c_{B(S)}e_S$. These specializations can be made simultaneously, because
the allocated blocks are disjoint. The determinant of this selected
$j$-minor is therefore a nonzero polynomial in a bounded number of input
variables, and at a point with coordinates in $\{0,1\}$ has absolute
value at least $c_0^j$.

Its degree is at most $j(h-1)$, and it depends on at most
$jh(q+1)$ variables; both numbers depend only on $h,q$. Its coefficients
thus include one of magnitude at least $c_0^j/C_{h,q}$, because evaluation
at a $\{0,1\}$ point is a sum of at most $C_{h,q}$ coefficients.
Lemma~\ref{cube:polynomial-sublevel} consequently gives
\[
 \mathbb P(|\det A_{\mathrm{ax}}|\leq\kappa)
 \leq C\kappa^{\alpha_{h,q}},
\]
for normalized cube-parameter laws on bounded polytopes of nonnegligible
volume; $C$ may depend polynomially on reciprocal coefficient bounds,
normalized density bounds and domain-volume bounds.  One applies
that lemma on a fixed containing box, then uses the density bound
or divides by the domain volume.  For $h>1$, an admissible exponent
is $\alpha_{h,q}=1/(jh(q+1)\,j(h-1))$; for $h=1$ the determinant
is a nonzero constant and the small sublevel sets are empty.
The determinant need not be large at an interior point of the
sampling polytope; its coefficient bound is sufficient.  For random
principal coefficients, condition on them first.  The bound is
uniform because their magnitudes lie between $c_0$ and $c_1$.

\paragraph{Several axes and the order of the choices.}
Suppose there are $D$ output axes with disjoint principal parameter
blocks, and write $d_\ell$ for their selected determinants.  Let
$C_*u^\alpha$ bound their individual sublevel probabilities, with
$\alpha>0$ depending only on the fixed degree and cube-order bounds.
Choose $\kappa\le1$ so that
\[
       C_*D(2\kappa)^\alpha\le\eta/2.
\]
The union of $\{|d_\ell|\le2\kappa\}$ has mass at most $\eta/2$.
If $H_0$ bounds the entries of the selected Jacobian blocks, their
inverse norms on $|d_\ell|\ge\kappa$ are at most
$C_{h,q}H_0^{j-1}/\kappa$.  The full selected ideal Jacobian is
block diagonal, so its inverse norm is the maximum of these block
norms.  No determinant of growing degree is used.

For the product laws under consideration, normalize coefficient
intervals to fixed boxes.  The scalar cube factors have dimension
$q+1$, a bounded number of faces, and pre-bounded densities.  Thus
there are $F_0=O_{h,q}(a)$ source faces, each with a width-$2r$
strip of probability at most $C_*r$, after enlarging $C_*$ by the
normalized domain bounds.  Choose $r\le1$ with
$C_*F_0r\le\eta/2$.  Smooth product cutoffs can vanish on the
width-$r$ boundary strips and on $|d_\ell|\le\kappa$, while
equaling one outside the width-$2r$ strips and $|d_\ell|\le2\kappa$.
For the latter factors use smooth functions of $d_\ell^2/\kappa^2$.
The resulting $\chi$ satisfies $\int(1-\chi)\rho\le\eta$ and
\[
 \|\nabla\chi\|_\infty
 \le C\left(F_0/r+
       \kappa^{-1}\sum_{\ell=1}^D\|\nabla d_\ell\|_\infty\right).
\]
The determinant derivative bounds have fixed degree.  Products of
cutoff factors cost the sum of their derivative bounds, since each
factor lies between zero and one.  Similarly, normalized product
coefficient densities have $L^1$ derivative bounds equal to the sums
of their one-factor bounds.  Hence
\[
 \|\chi\rho\|_{W^{1,1}}
 \le \|\rho\|_{W^{1,1}}+\|\nabla\chi\|_\infty
\]
has the claimed cost.  Apply Lemma~\ref{cube:image-perturbation}
and Remark~\ref{cube:perturbation-cost} to this cutoff.  In particular,
if $D$, the input dimension, and all logarithmic structural costs are
bounded by a fixed polynomial in a budget $p$, then specifying an error
$\exp(-p^{O(1)})$ requires only
\[
 \log\kappa^{-1},\quad \log B,
 \quad\log\delta^{-1},\quad\log\varepsilon^{-1}
       \leq p^{O(1)}.
\]
The order is: prescribe the discard error, choose the minor and
boundary cutoffs, choose the mollification scale, and only then
choose the size of the polynomial perturbation, also making its
first derivative smaller than half the reciprocal ideal inverse
bound.  These choices use no later lattice scale.  When the source
has small perturbation coefficients, they must first be written as
their width times unit-scale auxiliary variables; then the source
norms above are independent of that width and its smallness appears
in the image map itself.

These image estimates handle the continuous perturbation step.
Lemma~\ref{cube:product-block} supplies the Fourier truncation on
projection grids. The proof of Proposition~\ref{cube:comparison}
combines them with integer fiber counting and transfers Haar $L^1$
errors to ambient polynomial cube averages using
Corollary~\ref{sampling:normalization}. All these estimates use the
independent blocks and normalized domains specified above.

\subsection{Proof of the cube comparison}

\begin{proof}[Proof of Proposition~\ref{cube:comparison}]
We first retain the coefficient modes that factor through cube evaluation,
then count the spatial fibers and average the resulting jet law.
Finally we remove the small perturbations and expand the ideal law
into permitted site twists.

Discard the exceptional kernel tuples from
Lemma~\ref{cube:kernel-covers}. Fix a retained tuple and let $E_t$
be the evaluation map from polynomial coefficient arrays to Boolean
jets. Its real right inverse alone would not suffice for torus
characters: for example, $x\mapsto e(x)$ does not factor through
$x\mapsto2x$ on $\R/\Z$. The individual $M$-cover in that lemma
clears the denominators of the right inverse. Work on this cover,
pull the coefficient density $\mathcal D$ back to it, and normalize
against Haar probability. Write
\[
                       D_t=(E_t)_*\mathcal D.
\]
This notation means the density of the pushed-forward measure.
A Fourier character of the array that factors through the real jet
map factors on this cover as well: compose with a rational right
inverse having denominator dividing $M$.  Every other character is
removed by the assumed non-site-mode estimate, even with the
arbitrary site functions present.  Conditional Haar averaging
therefore leaves precisely $D_t$.

Apply Lemma~\ref{cube:fiber} to the spatial map from $(x,V)$ to the
root and differences of the spatial cube, one ambient coordinate at
a time.  Its normalized pivot uses the kernel columns and the root.
The resulting weight is a smooth real density multiplied by the
index and the image-lattice mask.  The original parent side lengths
can be chosen to make the fiber error as small as needed, including
the late sup bound of $D_t$.

Here is the spatial displacement estimate using the prescribed law.
For a fixed good kernel tuple, the normalized image of $(x,V_K)$,
where $K$ denotes the kernel columns, has a density $f_K$ with a pre
$C^1$ bound.  Indeed solve the output equations on the root and a
controlled kernel pivot, integrate the other smooth variables, and
differentiate the fixed normalized input density.  The nonkernel
support condition $|V_{ij}|\le c\sigma H_i/(mT_j)$ implies that its
contribution $Z_N$ to every normalized root or difference coordinate
has magnitude at most $C_q\sigma$.  Independence of the spatial
entries gives the full continuous density
\[
 f_t=\E_N f_K(\,\cdot-Z_N),\qquad
                 \|f_t-f_K\|_\infty\le C_{\rm pre}\sigma.
\]
The integer fiber still carries its full image index and membership
mask, even after this replacement of its real density.  These factors
are pre-bounded.  The resulting error in the functional is bounded by
$C_{\rm pre}\sigma$ times an ambient average of $D_t$.  Its Haar
mass is one, and Corollary~\ref{sampling:normalization}, with the
available late Fourier approximation, transfers that mass bound to
the ambient average.  Thus this error is pre-small without requiring
a pre supremum bound for $D_t$.  Narrow nonkernel spatial profiles
affect the earlier fiber error only through late powers of
$\sigma^{-1}$, which the parent side lengths absorb.

Write $\mathbf t_K$ for the fixed kernel tuple and $\mathbf t_N$
for all remaining parameter coordinates of the cube tuple. Condition
only on $\mathbf t_K$ and on $\mathbf t_N\bmod M$, not on the
values of $\mathbf t_N$.  The retained spatial weight $g_{\rm sp}$ is independent
of the unconditioned tuple entries.  To check this independence,
notice that the spatial fiber has already been summed over; the
parameter tuple is not conditioned on an equation $Vt=v$.
After the nonkernel spatial displacement is removed, the real
density uses only the fixed kernel tuple.  The image lattice
contains $M\Z^{q+1}$, so its membership and index are determined
by the fixed tuple residues modulo $M$.  The remaining cube
coordinates retain their product law across parameter axes,
including under the coordinatewise slice restrictions.
Up to an arbitrarily small error, the original functional is
the average over the retained kernel tuples and residue classes of
expressions of the form
\begin{equation}
 |U|^{-(q+1)}\sum_{\mathbf u\ {\rm cube\ in}\ U}
 g_{\rm sp}(\mathbf u)\prod_\omega P_\omega(u_\omega)
       \overline D\bigl(\chi(\mathbf u)\bigr),
 \qquad
 \overline D=\E[D_t\mid\mathbf t_K,\mathbf t_N\bmod M],
 \label{cube:ambient-functional}
\end{equation}
where $\chi(\mathbf u)$ is the collection of the Boolean jets of
$C_h(u_\omega)$ on the $M$-cover.  A finite Fourier expansion of the
spatial residue mask and Fourier smoothing of the normalized real
density express $g_{\rm sp}$ as a pre-cost mixture of products of
single-site factors.  The conversion from the normalization in
\eqref{cube:ambient-functional} to uniform genuine box cubes costs
at most an exponential in the ambient dimension.

It remains to approximate $\overline D$ in Haar $L^1$.  Use buffered
site charts so narrow that every alternating sum of lifts has
absolute components less than $1/2$.  Higher Boolean differences
vanish exactly: for $|S|>h$, the degree bound gives
$\Delta_SC_h=0$, while~\eqref{sampling:lift} gives
$\Delta_Sy_h=-\Delta_S\beta_h\in\Z^{d_h}$. The buffered chart
makes $|\Delta_Sy_h|_\infty<1/2$, so this integer vector is zero.
In these charts the lift jets are the actual real evaluations of their
coefficient arrays.  Reference Haar has Lebesgue measure in the
continuous split coordinates and grid volume $K_i^{-j}$ in a
projection coordinate, with the previously computed covolume
constants.  The remaining integral lattice coordinates on the cover
are independently uniform; their evaluated residues acquire only
the image index and membership mask modulo $M$.  Keep this mask
as an external multiplier of the output density; it does not
condition the independent principal coefficient laws.  More explicitly,
in the notation of Section~\ref{sampling:constrained-section}, the deck coordinates
of $\beta+l(z)$ are independent of the small coefficient lifts.
Since $l$ is integer-linear,
$E_tl(z_{\rm coeff})=l(E_tz_{\rm coeff})$; evaluated deck residues
therefore impose only the stated output mask and no additional
condition on the input projection coefficients.  This description
introduces no carries and no new equidistribution hypothesis.

For an enormous projection axis, replace integer coefficient
sampling by continuous sampling, keeping its index and residue
mask.  In Lemma~\ref{cube:fiber} the input scales are $K/T^\nu$,
the output scales are $K$, and the number of output rows is the fixed
number $j(h,q)$.  Here is why the normalized pivot has no hidden
$L$ loss.  On the kernel columns all $T_i=L$.  Divide the parameter
coordinates by $L$.  The retained real difference minor gives
linear polynomials, with pre-bounded coefficients in these
normalized coordinates, whose values are the individual Boolean
coordinates.  Their products, together with the constant, give a
pre-bounded real right inverse for the normalized degree-$h$ jet
matrix.  This uses only kernel monomials.  Cauchy--Binet then gives
a $j$-column minor with pre-bounded inverse: the number of possible
minors is polynomial in the number of input monomials, since $j$
is fixed.  The normalized entries are monomials in $t_i/T_i$ and
are bounded independently of $L$.  The coefficient density on these
coordinates may have widths of order $\sigma$; these affect its
cap and derivatives, or equivalently the inverse pivot after width
rescaling, only by allowed powers of $\sigma^{-1}$.

Since $T^\nu\le L^h$, the cell error is at most
$B_{\rm late}L^{hj}/(K/L^h)$.  After accounting for coefficient
widths and normalization, the error is thus bounded by
\[
                 B_{\rm late}\,L^{h(j+1)}/K.
\]
The logarithm of $B_{\rm late}$ is polynomial in the pre budgets
and $\log\sigma^{-1}$ and does not use $\log L$.  The exponent of
$L$ is independent of the number of coefficient variables and of
$J_0$.  Choose initially
$C_s>\max_{h,q}h(j(h,q)+1)$, over the fixed range of degrees and
cube orders.  Then the enormous cutoff $K>L^{C_s}$ and a later
choice of $A$ absorb the constant.  These real image
densities have late Lipschitz bounds, uniformly in the normalized
nonkernel tuple entries.  Continuous axes satisfy the same estimates
without a lattice error.  Normalized Riemann summation now makes
the dedicated long parameter tuples of the real and enormous axes
continuous.  Its error is bounded by a late constant times $M/L$,
including the boundary strips of their fixed-dimensional cube
polytopes.  Keep all parameter tuples and coefficients generating
inactive and moderate grid outputs discrete.

Remove the $\sigma$-perturbations in the real polynomial images.
First condition on all the variables generating those grid outputs.
Write every small coefficient as its prescribed center plus
$\sigma r$ in normalized coefficient coordinates.  The auxiliary
variable $r$ has a fixed smooth density with pre bounds; its density
does not have width $\sigma$.  Thus the common continuous input law
has pre derivative bounds.  Uniformly in the conditioned grid
variables, the true and ideal output maps $F_\sigma,F_0$ satisfy
\[
                 \|F_\sigma-F_0\|_{C^2}\le C_{\rm pre}\sigma.
\]
The normalized coordinates of the conditioned variables stay in a
fixed box, which gives this uniformity even for short inactive axes.

Apply Lemma~\ref{cube:image-perturbation} jointly to the real output
axes.  On each axis the disjoint principal blocks give the fixed-size
nonzero Jacobian minor constructed above.  The joint ideal Jacobian
is block diagonal.  Discard a pre-small sublevel set on each block and
a pre-small boundary region; a union bound over the axes controls the
discarded mass.  The individual inverse-minor and cutoff estimates
give the pre divergence bound in that lemma, and $C^1$ closeness
preserves invertibility along the interpolation.  Choose the discard
error and cutoffs, then the mollification scale, and only then
$\sigma$.  The resulting total variation error is pre-small,
uniformly under the conditioning.  Integrating the conditional
estimate proves a comparison of the joint real and grid law, not
only of its real marginal.  Grid summation on the enormous axes is
justified by the late Lipschitz bound for the perturbed density and
the pre bound for the ideal product-block density.  The index and
mask are pre-bounded and multiply absolute errors by at most a pre
factor.

After this replacement all axes factor.  Apply the product-block
estimate to the inactive and moderate grid laws and to the
continuous laws.  Every retained character
\eqref{cube:major-character} is a smooth character in $z/K$ times a
bounded-modulus character of $z$.  Boolean jets are linear
combinations of site data, so these characters split into products
of allowed site twists.  Buffered splitting-coordinate cutoffs
prevent periodic aliases and are equal to one on the ideal and true
supports.

There are two different error transfers in this argument.  The true
and ideal image densities both live on the same $M$-cover for the
fixed kernel tuple.  Their remaining discrete grid components have
$K\le L^{C_s}$; the enormous components have the smooth
interpolations just proved.  They therefore admit controlled late
Fourier approximations.  The absolute value of their difference has
the same type of Lipschitz control.  Corollary~\ref{sampling:normalization}
transfers its Haar integral, with the spatial weight, to the ambient
cube average in~\eqref{cube:ambient-functional}.  The arbitrary site
insertions are bounded by one in this estimate.  The final
ideal-to-truncated error is instead pointwise on the site charts,
so no joint cover for the many retained rational denominators is
needed.  Each retained term has only polynomially many such
denominators, whose product is pre-bounded.  Frequencies in the
bounded continuous ranges can finally be placed on a pre-fine mesh;
their characters change uniformly little on the bounded normalized
supports.  Averaging the corresponding coefficients makes the
mixture finite, with the same type of coefficient-mass bound.
\end{proof}

\subsection{Detection and finite-dimensional approximation}

\begin{lemma}[Detection on path-dependent slices]
\label{cube:pathwise-detection}
Fix the centers and a degree $r\ge0$, and assume the hypotheses and
conclusion of Proposition~\ref{cube:comparison} at cube order $r+1$.
Suppose a pre-large fraction of sampled maps detects a pre-bounded
function $\phi:U\to\C$ on permitted pre-cost slices against bounded
degree-$r$ local niltests of pre complexity.  The slices and niltests
may be chosen separately for each map.  Then $\phi$ has a pre-large
correlation on $U$ with a permitted site twist times a bounded
ambient degree-$r$ niltest of pre complexity.
\end{lemma}

\begin{proof}
Normalize the function and test caps first.
Let $\rho>0$ be a pre lower bound for the retained local correlations,
and let $\delta$ be a pre lower bound for the relative side counts of
the allowed slices.  Mesh their normalized endpoints finely enough
that passing to an inner mesh slice removes at most $\rho/4$ of its
points.  Its normalized correlation changes by at most $\rho/2$,
regardless of how rapidly the test oscillates.  On a sufficiently
long axis every permitted stride is at most $2/\delta$; enumerate
these strides, residues, and mesh endpoints.  An axis too short for
this comparison has pre-bounded length, so enumerate its resulting
subsets exactly.  If all axes are short, enumerate the actual slices
rather than their possibly redundant large-modulus descriptions.
The number of resulting slices is pre-bounded.  Restricting a local
niltest to the inner slice and reindexing changes its orbit
coefficients, which are unrestricted in Definition~\ref{prelim:niltest},
and preserves its degree and complexity.  Pigeonholing now gives one
common slice supporting a pre-large correlation for a pre-large
fraction of paths; their tests may still differ.

Lemma~\ref{prelim:dual-gowers} gives a pre-large $U^{r+1}$ norm on
this common reindexed box.  Its normalization is precisely
\eqref{prelim:box-gowers}; the overlap steps in its proof use the
cube-count weights in~\eqref{prelim:box-derivative}.  Put $q=r+1$.
The cube power is nonnegative on every path.  If the retained fraction is
$\tau$ and their local $U^q$ norm is at least $a$, then the average
local cube power is at least $\tau a^{2^q}$.  Apply
Proposition~\ref{cube:comparison} with error smaller than half this
number, taking $P_\omega=\mathcal C^{|\omega|}\phi$ at the vertices,
where $\mathcal C$ denotes complex conjugation.  If the coefficient
mass of its mixture is $M$, one of the mixed ambient integrals has
modulus at least $\tau a^{2^q}/(2M)$, by the triangle inequality;
the coefficients need not be positive.  Write
$f_\omega=\phi\,\mathcal C^{|\omega|}J_\omega$ and embed $U$ in
$H=\prod_i\Z/(4H_i)\Z$, extending these functions by zero.
If $\kappa=C_q(U)/|H|^{q+1}$, the mixed group average is $\kappa$
times the mixed box average, while
\eqref{prelim:ambient-gowers} multiplies each of the $2^q$ norms
by $\kappa^{1/2^q}$.  Gowers--Cauchy--Schwarz on $H$ therefore gives,
after exact cancellation of $\kappa$, the bound
\[
                  \prod_{\omega\in\{0,1\}^q}
          \|\phi\,\mathcal C^{|\omega|}J_\omega\|_{U^q(U)}
\]
for the modulus of that mixed integral.
The caps are pre-bounded, so at least one of these norms is
pre-large.  Divide that function by its cap and apply
Theorem~\ref{prelim:inverse} directly on the integer box $U$.
It supplies a degree-$r$ ambient niltest correlating with $\phi$
times the selected twist.  Moving that twist to the test side and
using closure under conjugation gives the claimed twisted niltest.
For $r=0$, a one-dimensional cube consists of two
independent sites; its integral is a product of two scalar means,
so a site twist alone detects $\phi$.  Restoring the original caps
has only pre cost.  This argument
does not require one common local niltest across paths; the common
slice and the lower bound for its cube norm are sufficient.
\end{proof}

We use the Hahn--Banach inverse-to-decomposition argument
of Gowers~\cite[Section~3.2]{GowersHahnBanach2010} to turn detection
into approximation. The version below also clips the dual witness
against a sampling density and makes the number of atoms independent
of the size of the ambient set. All pairings below are
normalized on a finite set, and complex convex hulls are absolutely
convex.

\begin{lemma}[Detection implies approximation]
\label{cube:separation}
Fix $\eta>0$.  Let $U$ be finite, $m:U\to[0,\infty)$ with $\E_Um\le C$, and
let $X$ be a seminorm satisfying
\[
                       X(g)\le C\E_Um|g|.
\]
Let $\mathcal J$ be a family of unit-bounded functions.  Suppose
that, for each specified cap $B$ and threshold $a>0$, every
$|v|\le B$ with $X(v)\ge a$ correlates by at least
$\delta(B,a)>0$ with some member of an allowed enlargement of
$\mathcal J$.  Assume also that a cap $C_0$ is fixed, and that the $L^1$ mass
of $m$ above $C_0$ is sufficiently small in terms of the target
$\eta$, $C$, $C_0$, and $\delta(CC_0/\eta,\eta/2)$.
For a family of sampling ensembles, this means $C_0$ is selected
first and the later ensemble scales are selected after this one
required tail accuracy.  It does not require a single fixed density
to have arbitrarily small tails above $C_0$.  Then every
$|f|\le1$ has an approximation by a bounded-coefficient-mass
combination of members of the enlarged $\mathcal J$ with
$X$-error at most $\eta$.  The coefficient mass and the number of
terms depend only on the quantitative data just stated.
\end{lemma}

\begin{proof}
Put $B=CC_0/\eta$ and $\delta_0=\delta(B,\eta/2)$, and fix the
single allowed enlargement for this cap and threshold, calling it
$\mathcal J$ in the proof.  It suffices to assume the explicit tail bound
\[
 \E m1_{m>C_0}\le\frac\eta C\min(1/8,\delta_0/4).
\]
Choose $M>\max(1,2/\delta_0)$.  Suppose $(9/8)f$ is outside
\[
 M\overline{\operatorname{aconv}}(\mathcal J)
                   +\eta\{g:X(g)\le1\}.
\]
The first summand is compact in this finite-dimensional space, so
the sum with the closed seminorm ball is closed.
Finite-dimensional separation, followed by normalization, gives a
functional $\chi$ such that
\[
 \Re\langle (9/8)f,\chi\rangle=1,\qquad
 X^*(\chi)\le\eta^{-1},\qquad
 \sup_{J\in\mathcal J}|\langle J,\chi\rangle|\le M^{-1}.
\]
The domination of $X$ implies pointwise
$|\chi|\le C m/\eta$: test the dual inequality on a function
supported at one point.  Clip $\chi$ along rays at $B$, obtaining
$v$.  The tail assumption gives
\[
 \E|\chi-v|\le\frac C\eta\E m1_{m>C_0}
                \le\min(1/8,\delta_0/4).
\]
Since $|f|\le1$, the normalization implies $\E|\chi|\ge8/9$.
Hence $\E|v|\ge8/9-1/8=55/72$, and radial clipping gives
\[
 X(v)\ge\eta\Re\langle v,\chi\rangle
       \ge\eta\E|v|^2\ge\eta(55/72)^2>\eta/2.
\]
Detection supplies $J\in\mathcal J$ with
$|\langle v,J\rangle|\ge\delta_0$.  Its cap is one, so
$|\langle\chi,J\rangle|\ge3\delta_0/4>M^{-1}$, a contradiction.
Dividing the resulting inclusion by $9/8$ gives
\[
 f\in (8M/9)\overline{\operatorname{aconv}}(\mathcal J)
                     +(8\eta/9)\{g:X(g)\le1\}.
\]

There is now a margin $\eta/9$ for a finite approximation.  First
approximate the closed-hull term by a finite combination with
coefficient mass at most $M$ and $X$-error at most $\eta/18$;
this is possible in finite dimension.  Then sample its atoms
independently, retaining their phases and total coefficient mass.
For $N$ samples the expected squared error in $L^2(1+m)$ is at most
$(1+C)M^2/N$, since the atoms have cap one.  Moreover
\[
 X(g)\le C(\E m)^{1/2}(\E m|g|^2)^{1/2}
                    \le C^{3/2}\|g\|_{L^2(1+m)}.
\]
A choice of $N$ depending only on $C,M,\eta$ gives a realization
with $X$-error at most $\eta/18$.  Its total error is at most
$8\eta/9+\eta/18+\eta/18=\eta$, and its number of terms is
independent of $|U|$.  The same sampling argument applies to a
bounded mixture.
\end{proof}

For the sampler take
$X(\phi)=\E_\psi\sup_{\text{slice,test}}
|\E_t\phi(\psi(t))\operatorname{test}(t)|$.
Every allowed slice has at least a pre fraction of the full
parameter mass, say $\delta_{\rm sl}$, and every test has been
normalized to cap one.  Before averaging over paths, its supremum
is at most $\delta_{\rm sl}^{-1}\E_{t\in\mathcal T}
|\phi(\psi(t))|$.  Consequently
\[
 X(\phi)\le\delta_{\rm sl}^{-1}\E_Um_U|\phi|,
                         \qquad \E_Um_U=1,
\]
where $m_U$ is the full one-site density.  No exchange of the
supremum with a conditional expectation is involved.  If $X(v)$ is
pre-large for a pre-bounded $v$, boundedness of this pathwise
supremum gives a pre-large fraction of detecting paths, to which
Lemma~\ref{cube:pathwise-detection} applies.
Corollary~\ref{sampling:normalization}
supplies the required tail above a fixed pre cap for $m_U$.
Choose its one required tail accuracy after the detector cap and
threshold in Lemma~\ref{cube:separation}, and then choose the later
ensemble scales.  That lemma gives the asserted approximation clause.
